\documentclass[12pt]{article}
\usepackage{fullpage}
\usepackage{amsmath,amssymb,amsthm}

\newtheorem{lemma}{Lemma}
\newtheorem{proposition}{Proposition}

\DeclareMathOperator{\Lk}{Lk}
\newcommand{\bbarotimes}{\mathbin{\overline{\otimes}}}
\newcommand{\cA}{\mathcal A}
\newcommand{\cH}{\mathcal H}
\newcommand{\C}{\mathbb C}

\begin{document}

\title{Proper proximality of irreducible graph-product von Neumann algebras}
\author{}
\date{}
\maketitle

\section*{Problem statement and interpretation}

I use the standard graph-product meaning of irreducible: the finite graph
\(\Gamma=(V,E)\) is not a non-trivial join, equivalently the complement
\(\Gamma^c\) is connected.  The traces are faithful normal traces.  For
\(U\subset V\), \(M_U\) denotes the graph product over the induced subgraph and
\(E_U:M\to M_U\) its canonical expectation.

Let \(A_N=\mathbb B(L^2N)\).  I write
\[
        \cA_N=(A_N^{\sharp_J})^*
\]
for the normal bidual of Ding--Kunnawalkam Elayavalli--Peterson (DKEP)
\cite{DKEP}.  If \(X\) is an \(N\)-boundary piece, \(K_X(N)\) is DKEP's
associated hereditary algebra and \(q_X\in\cA_N\) is the support projection of
\((K_X(N)^{\sharp_J})^*\).  Thus
\[
 \widetilde S_X(N)=
 \{T\in\cA_N:[T,a]\in q_X\cA_Nq_X\hbox{ for all }a\in JNJ\}.
\]
For \(X=\mathbb K(L^2N)\) write \(q_0\) and \(\widetilde S(N)\).  We use DKEP
Lemma 8.5 in the following form: for a separable finite factor, proper
proximality relative to \(X\) is equivalent to the non-existence of an
\(N\)-central state on \(\widetilde S_X(N)\) restricting to the trace.

\section*{Support and graph-product boundary corners}

Let \(p_{\rm nor}\in A_N^{**}\) be the central projection implementing the
simultaneous \(N\)- and \(JNJ\)-normal representation, so that
\(\cA_N\) is identified with \(p_{\rm nor}A_N^{**}p_{\rm nor}\).  If
\(B\subset A_N\) is a hereditary \(C^*\)-subalgebra which is invariant under
the multiplier actions of \(N\) and \(JNJ\), and \(s_B\in A_N^{**}\) is its
ordinary open support, then the paragraph before DKEP Remark 2.8 gives the
one-normal cutdown and DKEP Proposition 2.10 gives the simultaneous cutdown:
\(s_B\) commutes with \(p_{\rm nor}\), a contractive approximate unit of \(B\)
converges ultraweakly in \(\cA_N\) to
\[
        q(B)=p_{\rm nor}s_B=s_Bp_{\rm nor},
        \qquad
        (B^{\sharp_J})^*=q(B)\cA_Nq(B).                         \tag{1}
\]
For a DKEP boundary piece \(X\), Proposition 3.6 and the strong-bimodule
annihilator criterion of Proposition 2.6 imply that every
\(\sharp_J\)-functional vanishing on \(X\) vanishes on \(K_X^{\infty,1}(N)\),
hence on \(K_X(N)\); the converse follows from \(X\subset K_X(N)\).  Thus \(X\)
and \(K_X(N)\) have the same annihilator in \(A_N^{\sharp_J}\), and the normal
support \(q_X\) may be computed from the original boundary piece \(X\).

We shall repeatedly use two elementary support facts in \(A_N^{**}\).  First,
if \(d,c\) are projections and \(T^*T\in dA_N^{**}d\),
\(S^*S\in cA_N^{**}c\), then
\[
        T^*S\in dA_N^{**}c,                                      \tag{2}
\]
because \(T=Td\) and \(S=Sc\).  Secondly, if \(D,C\subset A_N\) are hereditary
with supports \(d,c\), \(e\) is a projection, and there are contractive
approximate units \(d_i\in D\), \(c_j\in C\) such that
\(d_ic_j\in eA_N^{**}c\) and \(c_jd_i\in cA_N^{**}e\) for all \(i,j\), then
\[
        d(cA_N^{**}c)d\subset eA_N^{**}e.                         \tag{3}
\]
Indeed, for \(x\in cA_N^{**}c\), the terms
\(d_ic_jxc_kd_l=(d_ic_j)x(c_kd_l)\) lie in \(eA_N^{**}e\), and successive
weak-* limits give \(dxd\).

For \(U\subset V\), put
\[
        \cH_U=L^2M\otimes_{M_U}L^2M
\]
as an \(M\)-\(M\) correspondence.  If \(\Omega_U=1\otimes1\), then
\(T_{\Omega_U}=e_{M_U}\).  Hence DKEP Example 5.11, together with Example 3.4,
identifies \(X_{M_U}\) with the coefficient boundary \(X_{\cH_U}\): one
inclusion follows because coefficients of \(M\)-\(M\) translates of
\(\Omega_U\) are the translated \(e_{M_U}\)-coefficients of Example 3.4, while
the other follows because \(X_{\cH_U}\) is hereditary, has \(M\) and \(JMJ\) in
its multiplier algebra, and contains \(e_{M_U}\), hence the whole corner
\(e_{M_U}A_Me_{M_U}\).  Let \(s_U\) be the ordinary support of \(X_{M_U}\).  By
(1),
\[
        q_U:=q_{X_{M_U}}=p_{\rm nor}s_U .                       \tag{4}
\]
If \(Z\subset U\), then \(e_{M_Z}\le e_{M_U}\); the preceding description gives
\(X_{M_Z}\subset X_{M_U}\), and hence \(s_Z\le s_U\).

We also need a harmless direct-sum observation.  If a correspondence
\(\cH=\oplus_i\cH_i\) has each \(X_{\cH_i}\) supported by a projection \(s\),
then every coefficient of every bounded vector in any Hilbert amplification of
\(\cH\) belongs to \(sA_M^{**}s\).  For finite-coordinate vectors this is (2)
applied to diagonal coefficients; for general bounded vectors, coordinate
projections preserve boundedness and the associated \(T\)-operators converge
strong-* to \(T_\xi\), so finite coefficients converge weak-*.

Finally, since \(X_{M_U}=X_{\cH_U}\) is hereditarily generated by positive
coefficients, it has a contractive approximate unit whose elements are norm
limits of finite sums of coefficients of finite fusions of \(\cH_U\) and
\(\overline{\cH_U}\).  This is obtained by applying continuous functions
vanishing at \(0\) to finite sums of positive coefficient generators; DKEP's
identities \(T_{\bar\xi}=T_\xi^*\) and \(T_{\xi\otimes\eta}=T_\eta T_\xi\)
convert the resulting monomials into such finite-fusion coefficients.  In the
tracial situation \(\overline{\cH_U}\cong\cH_U\).  By repeated use of the
CDHJKN fusion rule \cite[Theorem 5.4 and Proposition 5.5]{CDHJKN}, any finite
fusion using only labels contained in \(U\) decomposes as a direct sum of
\(\cH_Z\)'s with \(Z\subset U\); by the previous paragraph all its coefficients
are supported by \(s_U\).

\begin{lemma}[boundary-corner intersection]\label{lem:intersection}
For all \(U,W\subset V\),
\[
        q_W(q_U\cA_Mq_U)q_W\subset q_{U\cap W}\cA_Mq_{U\cap W}.   \tag{5}
\]
Consequently the palindromic product
\[
R_m=q_{U_m}q_{U_{m-1}}\cdots q_{U_1}q_{U_2}\cdots q_{U_{m-1}}q_{U_m}
\]
belongs to \(q_I\cA_Mq_I\), where \(I=\cap_iU_i\).
\end{lemma}

\begin{proof}
Let \(I=U\cap W\) and \(A=A_M\).  Choose the special approximate units
\(h_\alpha^W\in X_{M_W}\) and \(h_\beta^U\in X_{M_U}\) described above.  We
claim
\[
        h_\alpha^Wh_\beta^U\in s_IA^{**}s_U,\qquad
        h_\beta^Uh_\alpha^W\in s_UA^{**}s_I .                    \tag{6}
\]
It suffices, by norm approximation, to multiply two finite-fusion coefficients.
Let \(a=T_{\xi_1}^*T_{\xi_2}\) come from a finite fusion with all labels in
\(W\), and let \(b=T_{\eta_1}^*T_{\eta_2}\) come from a finite fusion with all
labels in \(U\).  DKEP's coefficient calculus gives
\[
        ab=T_{\xi_1\otimes\overline{\xi_2}\otimes\eta_1}^{\,*}
             T_{\eta_2}.                                      \tag{7}
\]
The left vector belongs to a finite fusion which, after first reducing the
\(W\)-part and then fusing with the \(U\)-part, decomposes by CDHJKN Proposition
5.5 into a direct sum of \(\cH_Z\)'s with \(Z\subset I\).  Its diagonal
coefficient is therefore supported by \(s_I\).  The diagonal coefficient of
\(T_{\eta_2}\) is supported by \(s_U\).  The rectangular support fact (2) gives
\(ab\in s_IA^{**}s_U\).  The adjoint argument gives the second inclusion in
(6).

Now apply (3) with \(D=X_{M_W}\), \(C=X_{M_U}\), and \(e=s_I\).  We obtain
\[
        s_W(s_UA^{**}s_U)s_W\subset s_IA^{**}s_I .
\]
Cutting this inclusion by \(p_{\rm nor}\) and using (4) gives (5).  The
palindromic assertion follows by induction from
\(R_{k+1}=q_{U_{k+1}}R_kq_{U_{k+1}}\).
\end{proof}

\section*{A finite boundary criterion}

\begin{lemma}\label{lem:finite}
Let \(N\) be a separable non-amenable finite factor.  Let \(X_1,\ldots,X_m\) be
boundary pieces and put \(q_i=q_{X_i}\).  Assume
\[
        R=q_mq_{m-1}\cdots q_1q_2\cdots q_{m-1}q_m\in q_0\cA_Nq_0. \tag{8}
\]
If \(N\) is properly proximal relative to every \(X_i\), then \(N\) is properly
proximal.
\end{lemma}

\begin{proof}
Assume \(N\) is not properly proximal.  By DKEP Lemma 8.5 there is an
\(N\)-central state \(\omega\) on \(\widetilde S(N)\) with
\(\omega|_N=\tau\).  Fix \(i\) and write \(q=q_i\).  Since \(q\) commutes with
\(N\) and \(JNJ\), \(q^\perp\widetilde S_{X_i}(N)q^\perp\subset\widetilde S(N)\).
If \(\omega(q^\perp)>0\), then
\[
        T\mapsto \omega(q^\perp Tq^\perp)/\omega(q^\perp)
\]
is an \(N\)-central state on \(\widetilde S_{X_i}(N)\).  On \(N_+\) it is
dominated by \(\omega(q^\perp)^{-1}\tau\), hence is normal; since \(N\) is a
factor its restriction is \(\tau\), contradicting relative proper proximality.
Therefore \(\omega(q_i^\perp)=0\) for all \(i\).

Extend \(\omega\) to a state on \(\cA_N\).  In its GNS representation every
\(q_i\) fixes the cyclic vector, so \(\omega(R)=1\).  Since \(R\) is a positive
contraction in \(q_0\cA_Nq_0\), \(\omega(q_0)=1\).  The u.c.p. map from
\(\mathbb B(L^2N)\) to \(q_0\cA_Nq_0\) used in DKEP's proof of Lemma 8.5 is
\(N\)-bimodular and is the identity on \(N\); composing it with \(\omega\) gives
a Connes hypertrace on \(\mathbb B(L^2N)\), contradicting non-amenability of
\(N\).
\end{proof}

\section*{Relative Bass--Serre paradox}

\begin{proposition}\label{prop:BS}
Let \(N=P*_{B}Q\) with trace-preserving expectations and suppose
\(Q=B\bbarotimes A\).  Assume that \(A\) contains a trace-zero unitary \(a\) and
that \(P\) contains unitaries \(w_1,w_2\) satisfying
\[
        E_B(w_i)=0,\qquad E_B(w_1^*w_2)=0 .
\]
Then \(N\) is properly proximal relative to \(X_B\).
\end{proposition}

\begin{proof}
Let \(H_P=L^2P\ominus L^2B\) and \(H_Q=L^2Q\ominus L^2B\).  In the AFP Fock
decomposition of \(L^2N\), let \(p\) and \(q\) be the projections onto the sums
of reduced words whose first letter lies in \(H_P\) and \(H_Q\), respectively;
then \(1=e_B+p+q\).

Write \(R_x=Jx^*J\).  Reduced-word inspection gives, for \(x\in P\ominus B\),
\[
        [p,R_x]=R_xe_B-e_BR_x,\qquad [q,R_x]=0,
\]
and for \(x\in Q\ominus B\) the analogous formula with \(p\) and \(q\)
interchanged; elements of \(B\) commute with \(p\) and \(q\).  The displayed
commutators are \(M\)- and \(JMJ\)-translates of \(e_B\), hence lie in the
\(X_B\)-mixing part.  DKEP Lemma 6.1, applied to the ultraweakly dense
\(*\)-algebra generated by \(P\) and \(Q\), gives \(p,q\in S_{X_B}(N)\).

Let \(\varphi\) be an \(N\)-central state on \(S_{X_B}(N)\) normal on \(N\).
Left multiplication by a centered \(P\)-letter sends \(qL^2N\) into \(pL^2N\);
moreover \(q\,w_1^*w_2\,q=0\) because \(E_B(w_1^*w_2)=0\).  Hence
\(w_1qw_1^*+w_2qw_2^*\le p\).  Similarly, since \(a\in Q\ominus B\), left
multiplication by \(a\) sends \(pL^2N\) into \(qL^2N\), so \(apa^*\le q\).
Centrality gives \(2\varphi(q)\le\varphi(p)\le\varphi(q)\), hence
\(\varphi(p)=\varphi(q)=0\).

The unitary \(h=w_1a\) is \(B\)-Haar: all non-zero powers are reduced
alternating words, so \(E_B(h^n)=0\) for \(n\ne0\).  Hence the projections
\(h^ne_Bh^{-n}\) are pairwise orthogonal and have equal \(\varphi\)-value; so
\(\varphi(e_B)=0\).  This contradicts
\(\varphi(1)=\varphi(e_B+p+q)=1\).  By DKEP Theorem 6.2, \(N\) is properly
proximal relative to \(X_B\).
\end{proof}

\section*{Application to graph products}

Fix \(v\in V\) and let \(L_v=\Lk_\Gamma(v)\).  The graph product splits as
\[
        M\cong M_{V\setminus\{v\}}*_{M_{L_v}}
        (M_{L_v}\bbarotimes M_v).                                \tag{9}
\]
Let \(C(v)=\{r\ne v:r\not\sim_\Gamma v\}\).  Since \(\Gamma^c\) is connected,
\(C(v)\ne\emptyset\).  If \(C(v)\) has distinct \(r,t\), take trace-zero
unitaries \(u_r\in M_r\), \(u_t\in M_t\) and put \(w_1=u_r,w_2=u_t\).  If
\(C(v)=\{r\}\), connectedness of \(\Gamma^c\) and \(|V|\ge3\) give
\(s\in L_v\) with \(s\not\sim_\Gamma r\); put \(w_1=u_r,w_2=u_su_r\).  In both
cases graph-product reduced-word calculus gives
\[
        E_{M_{L_v}}(w_i)=0,\qquad
        E_{M_{L_v}}(w_1^*w_2)=0 .
\]
The trace-zero unitary in \(M_v\) is the unitary \(a\) in the tensor factor of
(9).  Proposition \ref{prop:BS} shows that \(M\) is properly proximal relative
to \(X_{M_{L_v}}\) for every \(v\in V\).

Enumerate \(V=\{v_1,\ldots,v_m\}\).  Because no vertex belongs to its own link,
\(\cap_iL_{v_i}=\emptyset\).  Lemma \ref{lem:intersection} gives
\[
 q_{L_{v_m}}\cdots q_{L_{v_1}}q_{L_{v_2}}\cdots q_{L_{v_m}}
 \in q_{\emptyset}\cA_Mq_{\emptyset}.
\]
Here \(M_\emptyset=\C\), and \(X_{M_\emptyset}=\mathbb K(L^2M)\), since
\(e_{M_\emptyset}\) is rank one.

It remains to verify the hypotheses of Lemma \ref{lem:finite}.  The graph
product is separable.  By CDHJKN Main Theorem 0.5, factoriality can fail only
through a universal vertex of \(\Gamma\), or through a non-edge pair adjacent to
all other vertices with both vertex algebras two-dimensional.  These are,
respectively, an isolated vertex or an isolated edge in \(\Gamma^c\), impossible
because \(\Gamma^c\) is connected and has at least three vertices.  Thus \(M\)
is a factor.  By CDHJKN Proposition 6.3, amenability would force every non-edge
pair to be an isolated two-dimensional pair, again impossible for a connected
complement graph with at least three vertices.  Hence \(M\) is non-amenable.
Lemma \ref{lem:finite} applies, and \(M\) is properly proximal.

\begin{thebibliography}{9}

\bibitem{DKEP}
C.~Ding, S.~Kunnawalkam Elayavalli, and J.~Peterson,
\emph{Properly proximal von Neumann algebras},
Duke Math. J. \textbf{172} (2023), 2821--2894.

\bibitem{CDHJKN}
I.~Charlesworth, R.~de Santiago, B.~Hayes, D.~Jekel,
S.~Kunnawalkam Elayavalli, and B.~Nelson,
\emph{On the structure of graph product von Neumann algebras},
Publ. Res. Inst. Math. Sci. \textbf{61} (2025), 713--762.

\end{thebibliography}

\end{document}
